% ============================================================================
% HOUSE-PREAMBLE.tex -- the estate's house LaTeX style, as a file you copy in
% and \input rather than a block you copy out.
%
% See 00-RESEARCH-DOCUMENTATION.md, "The house preamble is a file, not a
% block". Copy this file to <repo>/docs/papers/house-preamble.tex (or
% docs/research/papers/ for a study paper) and \input it from every document
% in the repository that wants the house look. One copy per repository; every
% document in that repository \inputs the same one.
%
% WHY A FILE. The estate's first attempt at this shipped a template whose
% preamble was meant to be copied into each new document. Copying is what
% drifts: `ab-ove`'s two overview editions carried none of the seven house
% markers -- not the colors, the section formatting, the header, the status
% marker, the path macro, the geometry or the link style -- and nothing said
% so, because a convention about copying correctly has nothing to fail. The
% Beamer theme, which was always distributed as a file, propagated
% byte-identically over the same period. Distribute the file.
%
% ----------------------------------------------------------------------------
% THE CONTRACT. The including document defines these BEFORE \input, and
% carries its own \documentclass and \begin{document} -- this file has
% neither, so that one preamble can serve documents with different classes.
%
%   \paperstatus    Right-hand header, and the title block's status line.
%                   A study paper: DRAFT or VERIFIED, mirroring the study's
%                   own header. Any other document: its own status, named in
%                   that document's header comment as repurposed.
%   \headerleft     Left-hand header. Usually "Project --- Document title".
%   \pdftitleline   PDF metadata title.
%   \pdfauthorline  PDF metadata author.
%
% Optional, both defaulted below if the document does not define them:
%
%   \houselang      babel language for this edition, e.g. \def\houselang{polish}.
%                   Defaults to english. Set it in the edition file, not here:
%                   it is the one thing that is genuinely per-edition.
%   \editionsuffix  This edition's diagram-file suffix, e.g.
%                   \def\editionsuffix{.pl}. Defaults to empty. \dgm reads it,
%                   so a figure is called by slug and resolves per edition --
%                   see "Figures" below.
%
% A minimal English document:
%
%   \documentclass[11pt,a4paper]{article}
%   \newcommand{\paperstatus}{DRAFT}
%   \newcommand{\headerleft}{Project --- Overview}
%   \newcommand{\pdftitleline}{Project: an overview}
%   \newcommand{\pdfauthorline}{Author Name}
%   \input{house-preamble}
%   \begin{document}
%   ...
%
% Its Polish sibling differs by two lines, both before the \input:
%
%   \def\houselang{polish}
%   \def\editionsuffix{.pl}
%
% ----------------------------------------------------------------------------
% PACKAGES THIS FILE DOES NOT LOAD. The house set below is what every document
% needs. A document that needs more -- tikz, pifont, longtable, fancyvrb,
% amssymb -- loads it in its own preamble AFTER the \input, and says in a
% comment that it is an addition rather than part of the house set, so the
% next reader can tell the two apart at a glance.
% ============================================================================

% ----------------------------------------------------------------------------
% Encoding, then language. fontenc before inputenc before babel: babel reads
% the active font encoding when it sets up its shorthands, so loading it first
% gives an edition that is translated but typeset wrong.
% ----------------------------------------------------------------------------
\usepackage[T1]{fontenc}
\usepackage[utf8]{inputenc}

% The edition's language, defaulted so a monolingual document defines nothing.
% \PassOptionsToPackage rather than \usepackage[\houselang]{babel}: the option
% list here is a macro, and PassOptionsToPackage is the mechanism built for
% passing one.
\providecommand{\houselang}{english}
\PassOptionsToPackage{\houselang}{babel}
\usepackage{babel}

% ----------------------------------------------------------------------------
% The house package set.
% ----------------------------------------------------------------------------
\usepackage{geometry}
\usepackage{xcolor}
\usepackage{hyperref}
\usepackage{graphicx}
\usepackage{enumitem}
\usepackage{fancyhdr}
\usepackage{titlesec}
\usepackage{parskip}
\usepackage{booktabs}
\usepackage{array}
\usepackage{tabularx}

% ----------------------------------------------------------------------------
% Page geometry.
% ----------------------------------------------------------------------------
\geometry{
    a4paper,
    left=2.2cm,
    right=2.2cm,
    top=2.5cm,
    bottom=2.5cm
}

% ----------------------------------------------------------------------------
% House colors. Extracted from the estate's first formal documents
% (business_analysis.tex, pitch-deck-demium.tex) and unchanged since: a
% document from any repository should look like it came from the same shop.
% ----------------------------------------------------------------------------
\definecolor{primaryblue}{RGB}{41,128,185}
\definecolor{darkblue}{RGB}{31,78,121}
\definecolor{lightgray}{RGB}{236,240,241}
\definecolor{darkgray}{RGB}{52,73,94}
\definecolor{accentgreen}{RGB}{39,174,96}
\definecolor{accentorange}{RGB}{230,126,34}
\definecolor{accentred}{RGB}{231,76,60}

% ----------------------------------------------------------------------------
% Section formatting.
% ----------------------------------------------------------------------------
\titleformat{\section}
    {\color{primaryblue}\Large\bfseries\sffamily}
    {\thesection.}{0.5em}{}[\titlerule]

\titleformat{\subsection}
    {\color{darkblue}\large\bfseries\sffamily}
    {\thesubsection}{0.5em}{}

\titleformat{\subsubsection}
    {\color{darkgray}\normalsize\bfseries\sffamily}
    {\thesubsubsection}{0.5em}{}

% ----------------------------------------------------------------------------
% Header and footer. \paperstatus rides in the right-hand header of every
% page, so a draft cannot be mistaken for a final document at any page a
% reader happens to open it at -- which is the whole reason the marker is in
% the header rather than only on the title page.
% ----------------------------------------------------------------------------
\pagestyle{fancy}
\fancyhf{}
\renewcommand{\headrulewidth}{0.4pt}
\fancyhead[L]{\small\sffamily\color{darkgray} \headerleft}
\fancyhead[R]{\small\sffamily\color{accentred} \paperstatus\ --- \today}
\fancyfoot[C]{\thepage}

% ----------------------------------------------------------------------------
% Links and PDF metadata.
% ----------------------------------------------------------------------------
\hypersetup{
    colorlinks=true,
    linkcolor=primaryblue,
    urlcolor=primaryblue,
    pdfauthor={\pdfauthorline},
    pdftitle={\pdftitleline}
}

% ----------------------------------------------------------------------------
% Figures. A diagram's source is one .mmd file rendered to vector PDF by the
% repository's render step; the document includes the rendered file.
%
% \dgm resolves the path from a slug and the edition, so the body text names a
% diagram the way the diagrams directory does and never writes a path:
%
%   \includegraphics[width=0.78\linewidth]{\dgm{a1-system-context}}
%
% resolves to ../diagrams/rendered/a1-system-context.pdf in the default
% edition and ../diagrams/rendered/a1-system-context.pl.pdf in an edition that
% set \def\editionsuffix{.pl}. A repository whose editions share one set of
% diagrams simply never sets \editionsuffix, and every edition resolves to the
% same file.
%
% Hardcoding the path instead is what the estate did first, and it costs one
% edit per figure per edition: `ab-ove` maintained `b1-reader-loop.pdf` in the
% English file against `b1-reader-loop.pl.pdf` in the Polish one, by hand, in
% both. At three figures that is survivable; `agent-eval-bench` has 23.
%
% \housediagramdir is overridable for a document that does not sit one level
% below the diagrams directory -- define it before the \input.
% ----------------------------------------------------------------------------
\providecommand{\editionsuffix}{}
\providecommand{\housediagramdir}{../diagrams/rendered}
\newcommand{\dgm}[1]{\housediagramdir/#1\editionsuffix.pdf}

% ----------------------------------------------------------------------------
% House convenience commands.
%
% \repopath is for a path, a filename or an identifier copied from the code;
% it is deliberately the only one of these, so that a document does not grow a
% private synonym for it. \finding marks the sentence a section exists to
% deliver.
% ----------------------------------------------------------------------------
\newcommand{\repopath}[1]{\colorbox{lightgray}{\texttt{#1}}}
\newcommand{\finding}[1]{{\color{primaryblue}\textbf{#1}}}

% ----------------------------------------------------------------------------
% The house title block, as a command so that every document's first page
% looks the same and none of them reimplements it.
%
%   \housetitle{Title}{Subtitle or empty}{Author}{Affiliation or repo URL}{Provenance line}
%
% The provenance line is the one this estate keeps paying for the absence of:
% a document that presents another document names it here, so a reader who
% finds the two disagreeing knows which one is the source. For a study paper
% it names the companion study and the commit; for a project overview, the
% markdown it presents.
% ----------------------------------------------------------------------------
\newcommand{\housetitle}[5]{%
    \thispagestyle{plain}%
    \begin{center}
        {\LARGE\bfseries\sffamily\color{darkblue} #1}

        \ifx\relax#2\relax\else
            \vspace{0.3cm}

            {\Large\sffamily\color{darkgray} #2}
        \fi

        \vspace{0.6cm}

        {\large #3}

        \vspace{0.2cm}

        {\normalsize\color{darkgray} #4}

        \vspace{0.2cm}

        {\normalsize\color{accentred}\sffamily \paperstatus\ --- \today\ --- #5}
    \end{center}

    \vspace{0.4cm}
}
